\documentclass{../cls/ua-article}
% !TeX encoding = UTF-8
\UASetUniversity{Universidad Austral}
\UASetFaculty{Facultad de Ingeniería}
\UASetDepartment{Departamento de Computación}
\UASetCampus{Pilar}
\UASetCareer{Ingeniería Informática}
\UASetCommission{Comisión A}
\UASetProfessor{Equipo docente}
\UASetSemester{Segundo cuatrimestre 2026}
\UASetCourse{Sistemas Distribuidos}
\UASetDate{2026}

\UASetClassNumber{16}
\UASetClassTitle{Rúbrica de defensa final}
\UASetDocKind{Rúbrica formal de evaluación}
\title{Rúbrica formal de evaluación --- Defensa final}
\author{\UAProfessor}
\date{\UADate}
\begin{document}
\maketitle

\section{Filosofía de evaluación}
Se evalúa calidad de decisiones, no cantidad de tecnologías mencionadas.

\section{Dimensiones y ponderación}
\begin{center}
\begin{tabular}{p{5.2cm} p{7cm} c}
\toprule
Dimensión & Descripción & Peso \\
\midrule
Modelado del problema & ¿identifica restricciones y dominio? & 15 \\
Arquitectura global & ¿la propuesta es coherente y modular? & 15 \\
Consistencia y replicación & ¿la elección está justificada? & 15 \\
Manejo de fallas & ¿anticipa particiones, retries y degradación? & 15 \\
Seguridad & ¿modela amenazas y controles? & 10 \\
Dinámica real & ¿considera latencia, ancho de banda o costo físico? & 10 \\
Observabilidad & ¿puede operar y debuggear el sistema? & 10 \\
Trade-offs & ¿son explícitos, correctos y defendibles? & 10 \\
\bottomrule
\end{tabular}
\end{center}

\section{Escala cualitativa}
\begin{itemize}
\item Insuficiente: describe componentes sin comprender consecuencias.
\item Aceptable: usa conceptos correctos, pero de forma superficial.
\item Bueno: justifica decisiones con al menos un trade-off claro.
\item Excelente: anticipa fallas, límites y objeciones técnicas.
\item Nivel investigación: propone un insight fuerte o modelado original.
\end{itemize}

\begin{uakeybox}
Criterio de oro: el diseño debe poder sobrevivir, con sus supuestos explícitos, a una discusión técnica seria sobre producción real.
\end{uakeybox}
\end{document}
